\section{Conclusion}
\label{sec:conclusion}

This letter presented a computationally efficient framework for generating radar-oriented nonlinear sea surfaces across multiple environmental conditions. The framework begins with an Elfouhaily spectral realization, applies a nonlinear Lie-transform correction to the background surface, replaces a prescribed short-wave band with equal-energy directionally organized wind waves, and introduces localized curling geometry at selected steep crests. A facet-based two-scale scattering calculation and a statistical complex-echo model then connect the multiscale geometry to radar observables. This hierarchical construction retains explicit control over sea state, propagation direction, short-wave organization, and breaker morphology while avoiding the cost of resolving the complete hydrodynamic and full-wave electromagnetic evolution of every scene.

Validation was performed progressively from global statistics to local geometry, local scattering, full-scene observations, and computational practicality. The constrained nonlinear surfaces reduced the low-wind total-MSS RMSE from 0.00560 to 0.00223 relative to Cox--Munk and from 0.00523 to 0.00188 relative to the Gu\'erin reference, while preserving monotonic and saturating high-wind behavior. The conditioned curling waves reproduced the observation-constrained front-face-angle and front--rear-asymmetry ranges, and their unit-area radar-facing response produced a median enhancement of 5.069~dB, within the 2.25--7.85-dB literature-derived breaking-stage envelope. At the full-scene level, the proposed echoes improved the two-sample KS distance from 0.1254 to 0.0592, the Wasserstein distance from 0.1266 to 0.0380, and the logarithmic CCDF-tail RMSE from 1.7948 to 0.4420 relative to the linear baseline. The largest tested $1024\times1024$ grid required 0.8884~s of computation, and a stored $40\times40~\mathrm{m}^2$ realization required 0.0776~s on average, corresponding to an estimated 21.57~h for one million serial samples.

The present results support the proposed method as a controllable and scalable generator for large radar-oriented synthetic datasets, but they do not constitute a complete hydrodynamic or coherent full-wave description of breaking seas. The local electromagnetic experiment validates a facet-based response proxy rather than absolute breaker NRCS, and the accepted full-scene comparison does not insert localized curling events as a separate echo class. Foam, spray, entrained air, diffraction, multipath interference, and polarization-dependent breaker dynamics also remain outside the current formulation, while the simulated slow-time persistence is stronger than that of the measured record. Future work will therefore target held-out sea states and radar geometries, time-evolving breaker populations, selective higher-fidelity scattering for the most nonlinear regions, and downstream learning experiments that quantify how the generated diversity improves radar detection and recognition.
